\section{승인된 이력 추적과 사용 차단}
\label{sec:audit}

이 절은 앞 절에서 수락한 거래 이력을 따라 차단 대상을 찾고,
집행 결과를 확인하는 절차를 설명한다.
Fig.~\ref{fig:worked-flow}의 $A$에서 시작하면 생산·이전을 거친
$Z,E,W$가 남은 대상이고, 이미 발급된 $K$는 완료된 이력으로 남는다.
\begin{figure*}[t]
\centering
\begin{tikzpicture}[x=1cm,y=1cm,>=Stealth,
 obj/.style={draw,rounded corners,minimum width=10mm,minimum height=6mm,font=\small},
 ev/.style={draw,fill=gray!10,minimum width=14mm,minimum height=6mm,font=\scriptsize},
 target/.style={obj,very thick,fill=blue!5}]
 \node[obj] (a) at (0,0) {$A$};
 \node[ev] (s) at (1.6,0) {Split};
 \node[obj] (b) at (3.1,.9) {$B$};
 \node[obj] (c) at (3.1,-.9) {$C$};
 \node[ev] (t) at (4.7,.9) {Transfer};
 \node[obj] (v) at (6.3,1.6) {$V$};
 \node[target] (z) at (6.3,.1) {$Z$};
 \node[ev] (p) at (8,1.6) {Proceed};
 \node[obj] (d) at (9.6,1.6) {$D$};
 \node[obj] (x) at (9.6,.1) {$X$};
 \node[ev] (pr) at (11.3,1) {Process};
 \node[target] (e) at (13,1.7) {$E$};
 \node[target] (w) at (13,.2) {$W$};
 \node[ev] (is) at (4.7,-.9) {Issue};
 \node[obj] (k) at (6.3,-1.3) {$K$};
 \foreach \from/\to in {a/s,s/b,s/c,b/t,t/v,t/z,v/p,p/d,d/pr,pr/e,pr/w,c/is,is/k}
   \draw[->] (\from)--(\to);
 \draw[->,dashed] (x)--(pr);
 \node[font=\scriptsize,anchor=west,align=left] at (8,-1.1)
 {Targets: $Z,E,W$ (thick borders)\\$K$: terminal claim; $X$: other parent};
\end{tikzpicture}
\caption{Worked accounting flow after accepted production and transfer. The dashed edge is an actual input edge outside the forward traversal from $A$; the figure shows it only to explain the process accounting.}
\label{fig:worked-flow}
\end{figure*}

\subsection{선행·후속 추적}

추적은 확정된 하나의 원장 시점을 기준으로 한다.
원장 시작부터 그 시점까지의 기록 구간을 $H_t$라 쓰며,
확정 블록 해시 또는 그에 대응하는 수락 순서로 특정한다.
모든 조회는 같은 원장·시점에서 수행하고, 복호화에는 외부 승인 후
복원한 유효 감사키를 사용한다. 감사 기록의 전체 값 목록을 복호화한 뒤
거래 형식에 따라 부모 참조와 미래 식별자로 나눈다.
필요한 자료를 확인하지 못하면 추적을 실패로 처리하며,
이를 빈 경로나 미소비 상태로 간주하지 않는다.

\emph{선행 추적}(\Alg{TraceBackward})은 사용된 원료를 거슬러 올라가는 절차다.
객체의 생성 기록을 읽고, 그 기록에 암호화된 실제 부모 목록을 연다.
각 부모의 생성 기록으로 이동하는 일을 초기 등록까지 반복한다.
어떤 출력의 생성 기록을 읽었다고 해서 그 기록의 다른 출력까지
선행 경로에 포함하지는 않는다.

\emph{후속 추적}(\Alg{TraceForward})은 그 근거를 어디에 썼는지 따라가는 절차다.
객체의 생성 기록에서 미래 소비 식별자를 열고 소비 기록을 조회한다.
0이면 아직 쓰이지 않은 객체이므로 미소비 집합에 넣는다.
소비 이벤트가 있으면 타입·입력·생성 순서의 일치를 확인한다.
그 기록의 복호화된 부모 목록에도 현재 객체가 있는지 검사하고,
그 이벤트의 모든 출력으로 이어간다.
Claim과 출력 없는 종료는 종단으로 기록한다.

후속 추적은 분할·병합·가공에서 도달한 모든 출력을 포함한다.
폐기물이나 질량 0 출력도 뺄 수 없다.
다만 합류한 거래의 다른 입력을 거슬러 올라가지는 않는다.
그림의 $X$가 그 예다. 객체에서 시작하면 그 객체 하나,
기록에서 시작하면 그 기록의 전체 출력을 출발점으로 삼는다.
기록 시작의 선행·후속 그래프와 재방문 처리의 정확한 규칙은
부록~\ref{app:trace-details}에 명시한다.

결과에는 복원한 그래프와 미소비 회계 증표 $F_N$,
미해결 이전 증표 $F_V$, 종단 Claim과 종료 기록을 구분해 남긴다.
미소비 집합에는 타입 참조와 소비 식별자를 함께 보관한다.
이미 동결·사용 금지 상태인 객체도 아직 소비되지 않았다면 포함한다.

\subsection{차단 대상과 집행}

감사·동결 절차(\Alg{AuditAndFreeze})는 신뢰된 상태 관리자가
선택한 시작점의 차단을 승인한 뒤 실행한다.

전체 차단 완료를 목표로 하는 운영에서는 회계 객체의 소비와 차단
해제를 일시 중단하고, 그 중단이 적용된 뒤 확정된 $H_t$를 선택한다고
가정한다. 감사 조회와 동결 거래는 중단 중에도 허용한다.
중단은 아래의 공통 확인 시점 $H_r$까지 유지한다.
먼저 $H_t$에서 상태를 바꾸지 않고 후속 추적을 마친다.
추적에 성공한 뒤 동결 대상을
\begin{equation}
 T=F_N\cup F_V,\qquad t=(r,\nu)
 \label{eq:targets}
\end{equation}
로 고정한다. 각 $t$는 객체 참조 $r$와 소비 식별자 $\nu$의 쌍이다.
대상은 전체 미소비 집합과 정확히 같아야 한다.
이미 차단된 객체, 폐기물, 질량 0 객체와 집행에 실패한 객체도
대상 목록에 유지한다. Claim과 종료는 종단으로 남긴다.
대상이 비면 해당 시점에 차단할 미소비 객체가 없는 것으로 보고한다.

활성 대상에는 동결 거래를 제출하고, 이미 동결 또는 사용 금지된
대상은 그 차단 상태를 확인한다.
스마트 계약은 상태 관리자 권한, 타입별 식별자의 형식,
현재 미소비 여부와 허용 상태 전이를 검사한다.
감사자는 복호화한 결과에서 실제 객체의 식별자를 선택하고
전체 미소비 집합을 대상으로 유지할 책임을 맡는다.
대상별 오류 처리와 결과 형식은 부록~\ref{app:freeze-execution}에 둔다.

\subsection{차단 완료 판정}

집행을 마치면 전체 대상이 함께 차단되어 있는지 확인한다.
이를 위해 하나의 확정 시점 $H_r$에서 모든 대상을 다시 조회한다.
$H_r$은 $H_t$ 이후 또는 같은 시점이고, 같은 원장의 사용한 관찰과
거래 처리 결과를 포함해야 한다.

대상 $t$의 타입을 $\kappa_t$, 식별자를 $\nu_t$라 하자.
차단 상태 $\B$는 아직 소비되지 않았고 동결 또는 사용 금지인 경우다.
\begin{equation}
 \begin{aligned}
 \B_{H_r}(t)\ \Longleftrightarrow\quad
 &\spend^{\kappa_t}_{H_r}[\nu_t]=0\\
 {}\land\quad
 &\stat^{\kappa_t}_{H_r}[\nu_t]\in\{\stFrozen,\stRevoked\}.
 \end{aligned}
 \label{eq:blocked}
\end{equation}
\Alg{COMPLETE\_AT\_CHECKPOINT}는 다음 세 조건이 모두 맞을 때만 반환한다.
\begin{enumerate}
\item 후속 추적이 성공했고, 비어 있지 않은 전체 $T$가 대상 목록과
대상별 결과에 정확히 유지된다.
\item 위 조건에 맞는 공통 $H_r$을 확보했고,
제출 여부·처리 결과가 불명확한 거래가 없다.
\item $H_r$에서 모든 $t\in T$가 식~\eqref{eq:blocked}를 만족한다.
\end{enumerate}
정상 추적과 비어 있지 않은 $T$를 얻었지만 이 조건을 만족하지 못하면
\Alg{INCOMPLETE}로 보고하며 전체 대상과 대상별 결과를 보존한다.

완료는 $H_r$의 차단 상태에 관한 판정이며,
후속 소비는 차단이 유지되는 동안 거부된다.
전체 절차는 외부 감사와 여러 원장 거래를 연결해 수행한다.
